\clearpage\corrigesChapitre{3}{Préparer et doser une solution}
\corrigesSouples
\begin{blocLivre}{cor-3-1}\begin{corrige}{1}{Nommer les constituants}
\begin{questions}
\item Le glucose est le soluté ; l’eau est le solvant.
\item Non. Les entités de glucose sont dispersées dans la solution, elles ne cessent pas d’exister.
\end{questions}
\end{corrige}\end{blocLivre}
\begin{blocLivre}{cor-3-2}\begin{corrige}{2}{Concentration directe}
\begin{questions}
\item $V=\qty{0,200}{L}$.
\item $c_m=3,60/0,200=\qty{18,0}{g.L^{-1}}$.
\end{questions}
\end{corrige}\end{blocLivre}
\begin{blocLivre}{cor-3-3}\begin{corrige}{3}{Préparer par dissolution}
\schema[corrige]{dissolution}
\begin{questions}
\item $m=8,00\times0,2500=\qty{2,00}{g}$.
\item Dissoudre, transférer et rincer quantitativement dans la fiole de 250,0 mL, compléter au trait, boucher puis homogénéiser.
\end{questions}
\end{corrige}\end{blocLivre}
\begin{blocLivre}{cor-3-4}\begin{corrige}{4}{Une dilution par dix}
\schema[corrige]{dilution}
\begin{questions}
\item $F=100,0/10,0=10,0$.
\item $c_f=40,0/10,0=\qty{4,00}{g.L^{-1}}$.
\end{questions}
\end{corrige}\end{blocLivre}
\begin{blocLivre}{cor-3-5}\begin{corrige}{5}{Saturation}
\begin{questions}
\item $m_{\max}=200\times0,0500=\qty{10,0}{g}$.
\item 10,0 g se dissolvent au maximum ; 2,0 g restent solides si l’équilibre est atteint.
\end{questions}
\end{corrige}\end{blocLivre}
\begin{blocLivre}{cor-3-6}\begin{corrige}{6}{Choisir le prélèvement}
\begin{questions}
\item $V_p=c_fV_f/c_m=3,00\times100,0/15,0=\qty{20,0}{mL}$.
\item Une pipette jaugée de 20,0 mL et une fiole jaugée de 100,0 mL, avec propipette.
\end{questions}
\end{corrige}\end{blocLivre}
\begin{blocLivre}{cor-3-7}\begin{corrige}{7}{Lire l’étalonnage}
\schema[corrige]{etalonnage}
\begin{questions}
\item La lecture conduit à $c_m=\qty{60}{g.L^{-1}}$. La construction verte montre les projections.
\item Oui : 1,030 se situe dans l’intervalle de mesures de la gamme.
\end{questions}
\end{corrige}\end{blocLivre}
\begin{blocLivre}{cor-3-8}\begin{corrige}{8}{Deux masses différentes}
\begin{questions}
\item $c_m=5,00/0,1000=\qty{50,0}{g.L^{-1}}$.
\item $\rho=103,0/100,0=\qty{1,030}{g.mL^{-1}}$. Le premier quotient utilise seulement la masse de sel, le second toute la solution.
\end{questions}
\end{corrige}\end{blocLivre}
\clearpage
\begin{blocLivre}{cor-3-9}\begin{corrige}{9}{Deux dilutions successives}
\begin{questions}
\item Après la première dilution : 20,0 ; après la seconde : $\qty{4,00}{g.L^{-1}}$.
\item $F_{\mathrm{global}}=4\times5=20$. Les facteurs se multiplient.
\end{questions}
\end{corrige}\end{blocLivre}
\begin{blocLivre}{cor-3-10}\begin{corrige}{10}{Un trait dépassé}
\begin{questions}
\item Le volume est trop grand pour la même masse de soluté : la concentration est trop petite.
\item Non. Retirer une solution homogène retire eau et soluté dans la même proportion ; la concentration ne change pas. Il faut refaire la préparation ou la corriger par un calcul adapté.
\end{questions}
\end{corrige}\end{blocLivre}
\begin{blocLivre}{cor-3-11}\begin{corrige}{11}{Mélanger deux solutions}
\begin{questions}
\item $m=10,0\times0,100+4,00\times0,200=\qty{1,80}{g}$.
\item $c_f=1,80/0,300=\qty{6,00}{g.L^{-1}}$. Ce n’est pas la moyenne arithmétique des deux concentrations car les volumes diffèrent.
\end{questions}
\end{corrige}\end{blocLivre}
\begin{blocLivre}{cor-3-12}\begin{corrige}{12}{Une gamme insuffisante}
\demarche{Revenir à la solution initiale en utilisant le facteur de dilution.}
\begin{questions}
\item Elle place la solution mesurée dans le domaine couvert par les étalons, sans extrapolation hasardeuse.
\item $c_{\mathrm{initiale}}=5\times36,0=\qty{180}{g.L^{-1}}$.
\end{questions}
\retiens{La concentration mesurée est celle de la solution effectivement placée dans la gamme.}
\end{corrige}\end{blocLivre}

